\section{A stable correlation face and global lift size}\label{app:rank-one}
This appendix compares global formulation size with the local gap and
spatial resources in the main text. All matrix $\ell_1$ norms here are
\emph{entrywise} sums, and $\|H\|_\infty=\max_{ij}|H_{ij}|$ for a cost
matrix. Scalar inequalities count as nonnegative cone factors.

\subsection{The face and inherited exact lower bounds}
For $n\ge1$ define the compact sets
\[
 \mathcal R_n=\{W\ge0:\rank W\le1,\ We\le e,\ W^Te\le e\},
 \qquad K_n=\conv\mathcal R_n.
\]
For a nonzero generator, its row and column marginals $r,c$ and total
$S$ satisfy $W=rc^T/S$, $0\le r,c\le e$, and $e^Tr=e^Tc=S\le n$.
Indeed, writing $W=uv^T$ and computing its marginals proves this identity.
Then $\operatorname{tr}W=r^Tc/S\le1$. Equality forces
$r_i(1-c_i)=c_i(1-r_i)=0$ for every $i$, so $r=c=1_A$ for a nonempty
set $A$. Thus the trace-one face is the hull of $1_A1_A^T/|A|$.

\begin{proposition}[Explicit correlation face]\label{prop:correlation-face}
In $K_{2m}$ put
\[
 T(W)=e^TWe,\quad D(W)=1-\operatorname{tr}W,\quad
 B(W)=\sum_{i=1}^m W_{i,m+i},
\]
\[
 F=\{W\in K_{2m}:D=B=0,\ T=m\}.
\]
Then $F$ is a face affinely isomorphic to
$\operatorname{COR}(m)=\conv\{xx^T:x\in\{0,1\}^m\}$.
The forward map is $X=mW_{[m],[m]}$, and, with $x=\diag X$, its inverse is
\begin{equation}\label{eq:correlation-inverse}
 W=\frac1m\begin{pmatrix}
 X&xe^T-X\\ ex^T-X&ee^T-xe^T-ex^T+X
 \end{pmatrix}.
\end{equation}
Moreover, the affine functional
\[
 g(W)=m-T(W)+(2m+1)B(W)+(4m+1)D(W)
\]
satisfies $g\ge B+D\ge0$ on $K_{2m}$ and exposes exactly $F$.
\end{proposition}
\begin{proof}
Within the trace face, $B=0$ selects sets $A$ with at most one element
from each pair $\{i,m+i\}$. Their totals are $|A|\le m$; maximizing the
total selects exactly one element per pair. These successive faces give
atoms $(x,e-x)(x,e-x)^T/m$. Expanding their blocks gives
\eqref{eq:correlation-inverse}, which extends affinely to their hull.
For exposure, on a nonzero generator the inequality
$|a-b|\le a(1-b)+b(1-a)$ gives $\|r-c\|_1\le2SD$.
Also $r_i+r_{m+i}-r_ir_{m+i}\le1$, so
\[
 S\le m+\sum_i r_ir_{m+i}\le m+SB+\|r-c\|_1
       \le m+2mB+4mD.
\]
The zero generator satisfies the same affine inequality, hence so does
$K_{2m}$. Substitution proves $g\ge B+D$ and the asserted zero set.
\end{proof}

Adding affine equations to a lift and composing its output map leaves its
cone unchanged. Thus every lift of $K_{2m}$ gives one of
$\operatorname{COR}(m)$ over the same cone, without a strict-feasibility
assumption. The established correlation-polytope bounds therefore imply:
\begin{itemize}
\item For fixed $d$ and $m\ge d$, a lift over $(\mathbb S_+^d)^q$ has
\[
 q\ge(3^d-1)^{-(1-1/d)}(1-3^{-d})^{-m/d}.
\]
For $d=2$ the stronger bound is $q\ge7^{-1/2}(9/7)^{m/2}$
\citep[Theorem 1]{FawziParrilo2013}.
\item The total dimension of an exact second-order-cone lift is at least
$7^{-1/2}(9/7)^{m/2}$. A Lorentz cone
$\{(t,z):\|z\|_2\le t,\ z\in\R^k\}$ decomposes into $k-1$ three-dimensional
Lorentz cones for $k\ge2$, each isomorphic to $\mathbb S_+^2$;
one-dimensional nonnegative factors and the two-dimensional Lorentz cone
require at most their dimension in such blocks. Apply the preceding bound.
\item An unrestricted PSD lift has matrix order at least
$2^{\alpha m^{2/13}}$ for an absolute $\alpha>0$, for sufficiently large
$m$ \citep[Theorem 5.4]{LRS2015}. Several PSD blocks count by total order.
\end{itemize}
These are inherited conic lower bounds through the displayed face. They
assert neither nonrepresentability by finitely many second-order cones nor
an exponential number of Lorentz cones of unrestricted dimension.
For comparison, $K_2$ is already nonpolyhedral: the face where its first
row and column sums equal one is the hull of
$W(t)=(1,t)^T(1,t)/(1+t)$, $0\le t\le1$. With $q=W_{11}$ its lower-right
entry is $q+1/q-2$, a strictly convex function on $[1/2,1]$. Its infinitely
many graph points are extreme. Positive-error LP bounds below consequently
carry information that exact finite LP nonrepresentability alone does not.

\subsection{A linear error bound and a quantitative section transfer}
\begin{proposition}[Stability]\label{prop:correlation-stability}
For every $W\in K_{2m}$,
\begin{equation}\label{eq:face-error}
 \operatorname{dist}_1(W,F)\le(136m+10)g(W).
\end{equation}
If a convex set $Q$ satisfies $K_{2m}\subset Q\subset K_{2m}+\epsilon U_1$,
and $H_F=\{D=B=0,T=m\}$, then each $Y\in Q\cap H_F$ is within
$(184m+6)\epsilon$ of $F$. In particular, with $A_m=m(184m+6)$,
\begin{equation}\label{eq:face-outer-transfer}
 \operatorname{COR}(m)\subset\{mY_{[m],[m]}:Y\in Q\cap H_F\}
 \subset\operatorname{COR}(m)+A_m\epsilon U_1.
\end{equation}
\end{proposition}
\begin{proof}
For a nonzero generator $W=rc^T/S$, round $r$ to a binary vector $a$.
Since $\min(t,1-t)\le2t(1-t)$ and $\|r-c\|_1\le2SD$,
\[
 u:=\|r-a\|_1\le6SD,\qquad v:=\|c-a\|_1\le8SD.
\]
Indeed $\sum r_i(1-r_i)\le\sum r_i(1-c_i)+\|r-c\|_1\le3SD$.
Write $k=\sum a_i$, and let $h,z$ count pairs with two and zero ones.
Then $k=m+h-z$ and the product difference inequality on $[0,1]^2$ gives
$h\le SB+\|r-c\|_1+u\le SB+8SD$.
Change one coordinate in each full or empty pair to obtain $b$ with exactly
one one per pair. Thus
\[
 \|a-b\|_1=h+z=2h+m-k\le2SB+22SD+|m-S|.
\]
Set $V=bb^T/m\in F$ and $\widetilde b=(S/m)b$. Outer-product norm
factorization, using that $r,c,\widetilde b$ have total $S$, gives
\begin{align*}
 \|W-V\|_1
 &\le\|r-b\|_1+\|c-b\|_1+3|S-m|\\
 &\le4SB+58SD+5|S-m|.
\end{align*}
Here the first two terms bound the comparison with
$\widetilde b\widetilde b^T/S$, and its distance from $V$ is $|S-m|$.
Since $S\le2m$ and
$|S-m|\le g+(2m+1)B+(4m+1)D$, the last bound is at most
$(18m+5)B+(136m+5)D+5g\le(136m+10)g$.
For $W=0$ the distance to any paired atom is $m$, while $g(0)=5m+1$.
Taking finite convex combinations proves \eqref{eq:face-error}.

For the stronger transfer retain the bound
$8mB+116mD+5|S-m|$ on each generator, including zero. Exposure gives
$(S-m)_+\le2mB+4mD$. In a convex decomposition of $W$, therefore,
\[
 \sum_\ell\lambda_\ell|S_\ell-m|
 \le4mB(W)+8mD(W)+m-T(W).
\]
The corresponding combination $V\in F$ consequently obeys
\[
 \|W-V\|_1\le28mB(W)+156mD(W)+5[m-T(W)].
\]
For $Y\in Q\cap H_F$ choose $W\in K_{2m}$ with
$\|Y-W\|_1\le\epsilon$. The coefficient norms of $B,D,T$ give
$0\le B(W),D(W)\le\epsilon$ and $m-T(W)\le\epsilon$.
The triangle inequality yields $(184m+6)\epsilon$.
The block projection has norm at most $m$, proving the last assertion.
\end{proof}
An immediate exact-penalty consequence is that for
$\lambda>(136m+10)\|H\|_\infty$, minimizing
$\langle H,W\rangle+\lambda g(W)$ over $K_{2m}$ has all minimizers in $F$
and value $\min_F\langle H,V\rangle$. Compare $W$ with the point $V$
provided by \eqref{eq:face-error}; the excess is at least
$[\lambda-(136m+10)\|H\|_\infty]g(W)$.

\subsection{Fine LP and SDP approximations}
\begin{proposition}[Explicit accuracy scales]\label{prop:rank-one-approximation}
Let $K_{2m}\subset Q\subset K_{2m}+\epsilon U_1$ with $Q$ convex.
If $\epsilon\le1/A_m$, any LP lift of $Q$ has $2^{\Omega(m)}$
inequalities. If $\epsilon\le1/(2A_m)$, any PSD lift of $Q$ has matrix
order at least $\exp(c(m/\log m)^{2/13})$ for an absolute $c>0$ and all
sufficiently large $m$.
\end{proposition}
\begin{proof}
For the LP assertion apply \eqref{eq:face-outer-transfer}. The inequalities
$\langle2\diag(a)-aa^T,X\rangle\le1$, $a\in\{0,1\}^m$, have coefficient
norm at most one and slack $(1-a^Tx)^2$ at $xx^T$. The projected set is
therefore between $\operatorname{COR}(m)$ and the polyhedron defined by
these inequalities with right side two. The fixed-dilation lower bound of
\citet[Theorem 6(i), author version p.\ 18]{BraunEtAl2015} gives the result.
This particular sandwich has a small SDP relaxation and does not imply
the next assertion.

For that assertion use the sign correlation polytope
$\operatorname{SCOR}(m)=\conv\{yy^T:y\in\{-1,1\}^m\}$.
We first establish its fixed-error consequence of the pseudo-density
bounds in \citet[Theorem 3.8, equation (3.11), and Theorem 5.3]{LRS2015}.
For odd $k\ge3$ their construction supplies a degree-$k$ pseudo-density
$D$ on $\{0,1\}^k$ with
\[
 \E D=1,\quad\|D\|_\infty\le k^{3/2},\quad
 \E Df_k=-1/(4k^2),\qquad
 f_k(x)=\frac{(\sum x_i-k/2)^2-1/4}{k^2}.
\]
Here degree-$k$ means $\E Dp^2\ge0$ for $\deg p\le k/2$.
The cited quantitative theorem states that if $f:\{0,1\}^k\to[0,1]$
and $\E Df<-\delta$, then the matrix $M(S,x)=f(x_S)$, indexed by
$k$-subsets of $[m]$ and binary $x$, has, for $m\ge2k$, PSD rank at least
\begin{equation}\label{eq:lrs-input}
 \left[\frac{c_0\delta m}{k^3\|D\|_\infty\log m}\right]^{k/4}
 \left[\frac{\delta}{\|D\|_\infty}\right]^{3/2}\sqrt{\E f}.
\end{equation}
These pseudo-density and rank statements are established inputs.

Suppose $\operatorname{SCOR}(m)\subset R\subset
\operatorname{SCOR}(m)+\eta U_1$ with $\eta\le1/2$.
For every $k$-subset $S$ the affine function
\[
 L_S(Y)=\frac{\sum_{i,j\in S}Y_{ij}-1}{4k^2}
\]
is nonnegative on the sign polytope. Its coefficient norm is $1/(4k^2)$,
so $L_S+\theta\ge0$ on $R$, where $\theta=\eta/(4k^2)$.
At $y=2x-e$ its slack is $f_k(x_S)+\theta$. This shifted function is
in $[0,1]$, has mean at least $(k-1)/(4k^2)\ge1/(6k)$, and has
pseudo-expectation at most $-1/(8k^2)$. Taking $\delta=1/(16k^2)$
in \eqref{eq:lrs-input} gives
\begin{equation}\label{eq:shifted-psd-rank}
 \rank_{\rm psd}M\ge c_1k^{-23/4}
 \left[\frac{c_2m}{k^{13/2}\log m}\right]^{k/4}.
\end{equation}
If $R$ has a PSD lift of order $q$, these valid affine slacks have a PSD
factorization of order at most $q+1$. To justify this without a proper-lift
assumption, restrict the lift to the smallest PSD face containing its
feasible set. A finite combination of feasible matrices spanning all
feasible ranges is positive definite on this face. Conic duality for the
minimization of each nonnegative affine slack then writes that slack on
the affine slice as $\langle B,Y\rangle+\tau$, $B\succeq0$, $\tau\ge0$.
Appending the scalar factors $\tau$ and $1$ supplies the asserted
factorization. Choose odd $k$ comparable to
$a(m/\log m)^{2/13}$ with $a>0$ small enough that the bracket in
\eqref{eq:shifted-psd-rank} is at least 16. This proves the claimed
superpolynomial order for $R$.

Finally partition $W$ into four $m$-by-$m$ blocks and set
$J(W)=m(W_{11}-W_{12}-W_{21}+W_{22})$.
Its entrywise norm is at most $m$, and $J(F)=\operatorname{SCOR}(m)$.
The full-matrix transfer in Proposition~\ref{prop:correlation-stability}
shows that $R=J(Q\cap H_F)$ satisfies the preceding sandwich with
$\eta=A_m\epsilon\le1/2$. Affine sections and images preserve the lift
order, completing the proof.
\end{proof}
The scale in both assertions is explicit order $m^{-2}$. For compact
convex outer sets this norm guarantee is equivalent, by separation, to
uniform additive objective error $\epsilon$ over all costs with
$\|H\|_\infty\le1$. It does not cover only one objective, a restricted
cost family, or a larger accuracy tolerance. No spatial region-count
lower bound follows solely from these global lift-size bounds.

At the same SDP accuracy scale, a second-order-cone lift also needs
superpolynomial total cone dimension: decompose its Lorentz cones into
three-dimensional cones as above and embed them block-diagonally into
a PSD cone of order at most twice the original total dimension. This
transfers only the displayed superpolynomial approximate bound, whereas
the exact second-order-cone bound is exponential in $m$.
